\section{Variational Formulation}

In this section, we derive a weak formulation of the stochastic advection--diffusion problem introduced in 
Section~2. This formulation constitutes the basis for the well-posedness analysis and the finite element 
approximation developed later in the paper.

Let

\[
H=L^2(D)
\]

equipped with the inner product

\[
(y,v)_H
=
\int_D yv,dx,
\]

and the associated norm

\[
|y|_H
=
\left(
\int_D |y|^2,dx
\right)^{1/2}.
\]

We introduce the Sobolev space

\[
V=H_0^1(D),
\]

endowed with the norm

\[
|v|_V
=
|\nabla v|_{L^2(D)}.
\]

Since \(D\) is bounded, the Poincaré inequality implies that this norm is equivalent to the standard 
\(H^1(D)\)-norm.

The spaces \(V\), \(H\), and the dual space \(V'\) form the Gelfand triple

\[
V
\hookrightarrow
H
\hookrightarrow
V',
\]

where the embeddings are continuous and dense.

To derive the weak formulation, we multiply equation \eqref{eq:state} by a test function \(v\in V\) 
and integrate over the spatial domain (D). Using Green's formula together with the homogeneous Dirichlet boundary condition \eqref{eq:bc}, we obtain

\[
(dy,v)_H
+
a(y,v),dt
=
(u,v)_H,dt
+
(\sigma,dW(t),v)_H,
\]

where the bilinear form \(a:V\times V\rightarrow\mathbb{R}\) is defined by

\begin{equation}
a(y,v)
=
\kappa
\int_D
\nabla y\cdot\nabla v,dx
+
\int_D
(\mathbf b\cdot\nabla y),v,dx.
\label{eq:bilinear_form}
\end{equation}

The first term in \eqref{eq:bilinear_form} corresponds to diffusion, whereas the second term represents 
the advection mechanism.

Throughout this work, we assume that

\[
\mathbf b
\in
L^\infty(D)^d.
\]

Under this assumption, the bilinear form \(a(\cdot,\cdot)\) is continuous on \(V\times V\). Indeed, there 
exists a constant (C>0) such that

\[
|a(y,v)|
\le
C
|y|_V
|v|_V,
\qquad
\forall y,v\in V.
\]

Moreover, owing to the positivity of the diffusion coefficient \(\kappa\), the bilinear form satisfies a 
G{\aa}rding-type inequality

\begin{equation}
a(v,v)
\ge
\kappa
|\nabla v|_{L^2(D)}^2
C
|v|_{L^2(D)}^2,
\qquad
\forall v\in V,
\label{eq:garding}
\end{equation}

for some constant (C>0).

We are now in a position to define the notion of weak solution.

\begin{definition}
A stochastic process

\[
y:
[0,T]\times\Omega
\rightarrow
V
\]

is called a weak solution of problem
\eqref{eq:state}--\eqref{eq:ic}
if

\[
y
\in
L^2
\bigl(
\Omega;
L^2(0,T;V)
\bigr),
\]

\(y\) is \(\mathcal F_t\)-adapted, and for every test function \(v\in V\),

\begin{equation}
\begin{aligned}
(y(t),v)_H
&=
(y_0,v)_H
-
\int_0^t
a(y(s),v)\,ds
+
\int_0^t
(u(s),v)_H\,ds
\\
&\quad
+
\int_0^t
(\sigma\,dW(s),v)_H .
\end{aligned}
\label{eq:weak_form}
\end{equation}

holds almost surely for all \(t\in[0,T]\).
\end{definition}

The weak formulation \eqref{eq:weak_form} will serve as the starting point for the 
existence and uniqueness analysis presented in the next section.
